\documentclass[a4paper, serif, 10pt]{moderncv}

%% moderncv
% color and style
\moderncvcolor{black}
\moderncvstyle{banking}

% renew moderncv command to adapt for CJK colon
\renewcommand*{\cvitem}[3][.25em]{%
  \ifstrempty{#2}{}{\hintstyle{#2}：}{#3}%
  \par\addvspace{#1}}

\renewcommand*{\cvitemwithcomment}[4][.25em]{%
  \savebox{\cvitemwithcommentmainbox}{\ifstrempty{#2}{}{\hintstyle{#2}：}#3}%
  \setlength{\cvitemwithcommentmainlength}{\widthof{\usebox{\cvitemwithcommentmainbox}}}%
  \setlength{\cvitemwithcommentcommentlength}{\maincolumnwidth-\separatorcolumnwidth-\cvitemwithcommentmainlength}%
  \begin{minipage}[t]{\cvitemwithcommentmainlength}\usebox{\cvitemwithcommentmainbox}\end{minipage}%
  \hfill% fill of \separatorcolumnwidth
  \begin{minipage}[t]{\cvitemwithcommentcommentlength}\raggedleft\small\itshape#4\end{minipage}%
  \par\addvspace{#1}}

%% page layout/margins
\usepackage[top=2.5cm, bottom=2.5cm, left=1.5cm, right=1.5cm]{geometry}

\nopagenumbers{}

%% Babel config for English language
\usepackage[english]{babel}

%% fontspec
\usepackage{fontspec}

\IfFontExistsTF{Linux Libertine}{
  \setmainfont[Ligatures={TeX, Common}, Numbers=Lining, ItalicFont=Linux Libertine]{Linux Libertine}
}{}
\IfFontExistsTF{Linux Libertine O}{
  \setmainfont[Ligatures={TeX, Common}, Numbers=Lining, ItalicFont=Linux Libertine O]{Linux Libertine O}
}{}

%% CTeX
% CJK support, used to show CJK characters in the resume
%
% - heading=false: disable ctex heading
% - punct=kaiming: use kaiming punctuations styles for CJK
% - scheme=plain: use plain scheme, do not override \normalsize font size
% - space=auto: space settings for CJK characters
%
% CJK fonts are resolved per locale language, e.g. Noto Serif/Sans CJK JP
% for Japanese resumes, with fallback to the next region if not installed
%
% ref:
% - http://ctan.mirrorcatalogs.com/language/chinese/ctex/ctex.pdf
\usepackage[UTF8, heading=false, punct=kaiming, scheme=plain, space=auto]{ctex}

\IfFontExistsTF{Noto Serif CJK HK}{
  \setCJKmainfont{Noto Serif CJK HK}
}{
\IfFontExistsTF{Noto Serif CJK TC}{
  \setCJKmainfont{Noto Serif CJK TC}
}{}}
\IfFontExistsTF{Noto Sans CJK HK}{
  \setCJKsansfont{Noto Sans CJK HK}
}{
\IfFontExistsTF{Noto Sans CJK TC}{
  \setCJKsansfont{Noto Sans CJK TC}
}{}}

%% line spacing
\usepackage{setspace}
\setstretch{1.125}

%% URL styling - use normal text instead of monospace
\urlstyle{same}

%% Patch \httplink and \httpslink to handle full URLs
% The original moderncv macros always prepend http:// or https:// to URLs.
% This causes issues when the URL already has a protocol (e.g., https://example.com)
% resulting in malformed URLs like https://https://example.com.
% This patch checks if the URL already contains "://" and uses it directly if so.
% Note: We use \str_set:Nx (x-expansion) to expand macros like \@homepage first.
\ExplSyntaxOn
\str_new:N \g_yamlresume_url_str

\RenewDocumentCommand{\httpslink}{O{}m}{
  \str_set:Nx \g_yamlresume_url_str {#2}
  \str_if_in:NnTF \g_yamlresume_url_str {://}
    { \url{#2} }
    { \url{https://#2} }
}

\RenewDocumentCommand{\httplink}{O{}m}{
  \str_set:Nx \g_yamlresume_url_str {#2}
  \str_if_in:NnTF \g_yamlresume_url_str {://}
    { \url{#2} }
    { \url{http://#2} }
}
\ExplSyntaxOff

\name{陳嘉雯}{}
\title{資深營運經理 | 零售及供應鏈轉型}
\phone[mobile]{+852 9123 4567}
\email{ka.man.chan@example.com}
\homepage{https://www.kamanchan.com}

\address{中國香港，九龍，香港，香港九龍尖沙咀彌敦道100號}{}{}

\extrainfo{{\small \faLinkedin}\ \href{https://www.linkedin.com/in/chan-ka-man-ops}{@chan-ka-man-ops} {} {} {} • {} {} {} 
{\small \faMedium}\ \href{https://medium.com/@kamanchan-ops}{@kamanchan-ops}}

\begin{document}

\maketitle

\section{簡介}

\cvline{}{具備超過十年零售及供應鏈營運管理經驗，專注以數據驅動流程改善、成本控制及團隊效能提升。
%
\begin{itemize}
\item 曾管理逾80人跨區域團隊，負責年度營業額港幣2.5億元業務之損益管理
\item 擅長將營運指標可視化，透過Power BI及SAP推動補貨、庫存及門店流程優化
\item 主導六西格瑪及精實項目，三年內累計節省營運成本逾港幣1,200萬元
\item 具備跨部門協作與變革管理經驗，能協調採購、物流、銷售及客服團隊達成目標
\end{itemize}}
\section{教育背景}

\cventry{2018年9月–2020年6月}
        {碩士，供應鏈與物流管理，成績：3.8}
        {香港中文大學}
        {\url{https://www.cuhk.edu.hk}}
        {}
        {\begin{itemize}
\item 以系統動力學及數據模型研究零售補貨策略，畢業項目獲教授推薦作教學案例
\item 協助本地物流企業分析最後一公里配送成本，提出路線整合建議
\end{itemize}
\textbf{課程}：策略採購、物流網絡設計、營運分析、供應鏈風險管理、服務營運管理}

\cventry{2010年9月–2014年6月}
        {學士，工商管理（營運管理），成績：3.6}
        {香港理工大學}
        {\url{https://www.polyu.edu.hk}}
        {}
        {\begin{itemize}
\item 主修營運管理，建立流程分析、品質管理及成本控制基礎
\item 曾代表學系參與企業個案比賽，負責營運改善方案設計
\end{itemize}
\textbf{課程}：管理會計、作業管理、統計學、商業流程分析、項目管理}
\section{工作經歷}

\cventry{2021年3月至今}
        {營運經理}
        {恒匯零售集團有限公司}
        {\url{https://www.hanwui-retail.com.hk}}
        {}
        {\begin{itemize}
\item 管理3個區域、12間門店及80人團隊，負責年度營業額港幣2.5億元之營運及損益
\item 建立每週營運儀表板，追蹤銷售、客流、庫存周轉及人力成本，帶動門店生產力提升18\%
\item 主導智能補貨及中央庫存調撥項目，將缺貨率由9\%降至3.5\%，庫存周轉天數縮短22\%
\item 與採購、物流及市場部協作推出季節性推廣，連續兩年超額完成營業額目標
\item 制定標準作業程序及培訓計劃，將前線員工流失率降低15\%
\end{itemize}
\textbf{關鍵字}：零售營運、損益管理、流程改善、團隊領導、數據分析}

\cventry{2017年7月–2021年2月}
        {高級營運主任}
        {華駿物流有限公司}
        {\url{https://www.wahchun-logistics.com.hk}}
        {}
        {\begin{itemize}
\item 負責香港及華南地區第三方物流營運，管理倉儲面積逾10萬平方呎及60人團隊
\item 優化倉庫布局與揀貨路線，將訂單履行時間縮短25\%，準時交付率提升至98.5\%
\item 導入ISO 9001品質管理流程，完成內部審核及糾正預防措施，年度審核零重大缺失
\item 分析運輸及人力成本，重新議定承運商合約，每年節省港幣約380萬元
\end{itemize}
\textbf{關鍵字}：第三方物流、倉儲管理、ISO 9001、成本控制、準時交付}

\cventry{2014年7月–2017年6月}
        {營運分析師}
        {駿達供應鏈管理有限公司}
        {\url{https://www.chuntat-scm.com.hk}}
        {}
        {\begin{itemize}
\item 為零售客戶建立需求預測及補貨模型，將預測準確度提升12\%
\item 編製營運月報及管理層簡報，追蹤KPI、成本差異及改善措施進度
\item 參與供應商績效評核及訂單履行流程重組，縮短採購前置時間8\%
\end{itemize}
\textbf{關鍵字}：需求預測、供應商管理、KPI、報表分析}
\section{語言}

\cvline{粵語}{母語或雙語水平 \hfill \textbf{關鍵字}：母語、粵語}
\cvline{英語}{完全專業水平 \hfill \textbf{關鍵字}：HKDSE English Level 5、商務書信}
\cvline{普通話}{完全專業水平 \hfill \textbf{關鍵字}：普通話水平測試二級甲等、跨地區溝通}
\section{專業技能}

\cvline{營運管理}{專家 \hfill \textbf{關鍵字}：策略規劃、流程改善、損益分析、資源配置、風險管理}
\cvline{供應鏈與物流}{高級 \hfill \textbf{關鍵字}：採購、庫存控制、倉儲管理、第三方物流、需求預測}
\cvline{數據分析與工具}{高級 \hfill \textbf{關鍵字}：Excel、Power BI、Tableau、SQL、SAP}
\cvline{領導與溝通}{專家 \hfill \textbf{關鍵字}：團隊培訓、跨部門協作、變革管理、績效管理}
\section{榮譽}

\cventry{2023年11月}
        {香港零售管理協會}
        {傑出營運改善大獎}
        {}
        {}
        {表揚智能補貨與跨店庫存調撥項目，成功降低缺貨率並提升門店銷售表現。}

\cventry{2022年3月}
        {恒匯零售集團有限公司}
        {年度優秀團隊領袖}
        {}
        {}
        {帶領營運團隊連續兩年超額完成營業額及成本控制目標，獲管理層嘉許。}
\section{證書}

\cventry{2019年5月}
        {美國品質學會（ASQ）}
        {六西格瑪綠帶認證}
        {\url{https://asq.org/cert/six-sigma-green-belt}}
        {}
        {}

\cventry{2020年9月}
        {香港品質保證局}
        {ISO 9001:2015 內部稽核員證書}
        {\url{https://www.hkqaa.org}}
        {}
        {}

\cventry{2021年6月}
        {香港生產力促進局}
        {精實供應鏈管理證書}
        {\url{https://www.hkpc.org}}
        {}
        {}
\section{出版刊物}

\cventry{2022年8月}
        {零售營運數碼轉型實務指南}
        {香港物流協會月刊}
        {\url{https://www.hkla.org.hk/publications/retail-ops-digital-guide}}
        {}
        {\begin{itemize}
\item 分析零售企業如何以數據儀表板、需求預測及自動補貨提升營運效率
\item 提出由試點到規模化的數碼轉型路線圖，並列舉成本與效益評估方法
\end{itemize}}
\section{項目}

\cventry{2022年1月–2022年12月}
        {為12間門店建立自動補貨建議及跨店庫存調撥機制。}
        {智能補貨與庫存優化計劃}
        {\url{https://www.hanwui-retail.com.hk/projects/smart-replenishment}}
        {}
        {\begin{itemize}
\item 整合POS、庫存及客流數據，建立每週補貨建議模型
\item 與採購及物流團隊制定跨店調撥規則，減少滯銷及缺貨
\item 項目完成後缺貨率由9\%降至3.5\%，庫存周轉天數縮短22\%
\end{itemize}
\textbf{關鍵字}：自動補貨、庫存優化、數據儀表板}

\cventry{2019年3月–2019年12月}
        {以精實方法重整揀貨及包裝流程，提升第三方物流倉庫效率。}
        {倉庫流程精實改善項目}
        {\url{https://www.wahchun-logistics.com.hk/projects/lean-warehouse}}
        {}
        {\begin{itemize}
\item 進行時間研究及流程盤點，識別揀貨路線重複與包裝瓶頸
\item 重新劃分倉儲區域並導入標準作業程序，訂單履行時間縮短25\%
\item 建立每日改善會議及視覺化管理看板，前線參與度顯著提升
\end{itemize}
\textbf{關鍵字}：精實管理、倉儲流程、標準作業}
\section{興趣愛好}

\cvline{戶外運動}{行山、跑步、單車}
\cvline{閱讀}{營運管理、供應鏈、歷史}
\cvline{義工服務}{職涯指導、社區活動}
\section{志願服務}

\cventry{2019年1月至今}
        {職涯導師}
        {香港青年協會}
        {\url{https://www.hkfyg.org.hk}}
        {}
        {\begin{itemize}
\item 為準備投入職場的青年提供履歷撰寫、面試技巧及職涯規劃指導
\item 協助舉辦模擬面試及行業分享會，累計服務超過120名青年
\end{itemize}}
    
\end{document}